\begin{titlepage}
	
	% Defines a new command for the horizontal lines, change thickness here
	\newcommand{\HRule}{\rule{\linewidth}{0.5mm}} 
	
	% Center everything
	\center
	
	%----------------------------------------------------------------------------------------
	%	HEADING SECTIONS
	%----------------------------------------------------------------------------------------
	
	% Name of your university/college
	\textsc{\LARGE Università degli Studi di Padova}\\[1cm] 
	
	%----------------------------------------------------------------------------------------
	%	LOGO SECTION
	%----------------------------------------------------------------------------------------
	
	\includegraphics[height=5cm]{img/UniPd.png}\\[1cm]
	
	\includegraphics[height=1.5cm, width = 9cm]{img/MathDip.png}\\
	\textsc{\dipartimento}\\[1.2cm]
	% Major heading such as course name
	\textsc{\Large Scuola di Scienze}\\[0.5cm] 
	
	% Minor heading such as course title
	\textsc{\large Corso di Laurea in Informatica}\\[0.5cm] 
	
	\vspace{2.5cm}
	
	%----------------------------------------------------------------------------------------
	%	TITLE SECTION
	%----------------------------------------------------------------------------------------
	
	\HRule \\[0.4cm]
	{ \huge \bfseries PROGETTO DI WEB INFORMATION MANAGEMENT}\\[0.6cm] % Title of your document
	{ \huge \emph{Analisi di usabilità di un sito web} }
	\HRule \\[1cm]
	
	%----------------------------------------------------------------------------------------
	%	AUTHOR SECTION
	%----------------------------------------------------------------------------------------
	
	\vspace{1cm}
	
	\begin{minipage}{0.5\textwidth}
		\begin{flushleft} 
			\center
			\emph{\bfseries Autore: }\nomeStudente\ \cognomeStudente\ \\
			\emph{\bfseries Matricola: }\matricolaStudente \\
			\emph{\bfseries Sito analizzato: } \textit{http://www.animeforce.org/} \\
			\emph{\bfseries Periodo di analisi: } Marzo - Aprile 2019
		\end{flushleft}
	\end{minipage}
	
\end{titlepage}

\tableofcontents
\newpage